\global\long\def\defeq{\stackrel{\mathrm{{\scriptscriptstyle def}}}{=}}%
\global\long\def\norm#1{\left\Vert #1\right\Vert }%
\global\long\def\R{\mathbb{R}}%
\global\long\def\Diag{\mathrm{Diag}}%
\global\long\def\Rn{\mathbb{R}^{n}}%
\global\long\def\tr{\mathrm{Tr}}%
\global\long\def\diag{\mathrm{diag}}%
\global\long\def\rank{\mathrm{rank}}%
\global\long\def\ma{A}%
\global\long\def\mx{X}%
\global\long\def\ms{S}%
\global\long\def\mproj{P}%
\global\long\def\mi{I}%
\global\long\def\mzero{0}%
\global\long\def\oma{\overline{\ma}}%
\global\long\def\md{D}%
\global\long\def\mw{W}%
\global\long\def\omx{\overline{X}}%
\global\long\def\oms{\overline{S}}%
\global\long\def\ox{\overline{x}}%
\global\long\def\os{\overline{s}}%
\global\long\def\oy{\overline{y}}%
\global\long\def\ttag#1{\tag{#1}}%
\global\long\def\PriR{\mathcal{P}}%
 
\global\long\def\DualR{\mathcal{D}}%
 

\section{Interior Point Method for Linear Programs\label{sec:IPM}}

In this section, we study interior point methods. In practice, this
method reduces the problem of optimizing a convex function to solving
a small number of linear systems (``typically less than $30$'').
In theory, it reduces the problem to $\tilde{O}(\sqrt{n})$ linear
systems. 

We start by describing interior point method for linear programs.
We establish a polynomial bound and then later discuss implementation
details.

\section{Introduction}

Consider the primal linear program
\begin{equation}
\min_{\ma x=b,x\in\R^{n}_{\geq0}}c^{\top}x\ttag P\label{eq:tut_primal_LP}
\end{equation}
and its dual
\begin{equation}
\max_{\ma^{\top}y+s=c,s\in\R^{n}_{\geq0}}b^{\top}y\ttag D\label{eq:tut_dual_LP}
\end{equation}
where $\ma\in\R^{d\times n}$, $\R_{\geq0}$ denotes the set of nonnegative
real numbers, and vector inequalities are coordinatewise. The feasible
regions for the two programs are 
\[
\PriR=\{x\in\R^{n}_{\geq0}:\ma x=b\}\mbox{ and }\DualR=\{s\in\R^{n}_{\geq0}:\ma^{\top}y+s=c\text{ for some }y\}.
\]
We define their interiors: 
\[
\PriR^{\circ}=\{x\in\R^{n}_{>0}:\ma x=b\}\mbox{ and \ensuremath{\DualR^{\circ}}}=\{s\in\R^{n}_{>0}:\ma^{\top}y+s=c\text{ for some }y\}.
\]
One key property of the primal and dual linear programs is that the
dual objective $b^{\top}y$ is a lower bound on the primal objective
$c^{\top}x$ and that $x^{\top}s$ measures the gap between the two
objective values.

To motivate the main idea of the interior point method, we recall
the optimality condition for linear programs.
\begin{lem}[Duality Gap]
\label{lem:duality_gap}For any $x\in\PriR$ and any dual feasible
pair $(y,s)$ with $\ma^{\top}y+s=c$, the duality gap satisfies $c^{\top}x-b^{\top}y=x^{\top}s$.
In particular, we have
\[
\max_{\ma^{\top}\oy+\os=c,\os\in\R^{n}_{\geq0}}b^{\top}\oy-x^{\top}s\leq b^{\top}y\leq c^{\top}x\leq\min_{\ma\ox=b,\ox\in\R^{n}_{\geq0}}c^{\top}\ox+x^{\top}s.
\]
\end{lem}

\begin{proof}
Using $\ma x=b$ and $\ma^{\top}y+s=c$, we can compute the duality
gap as follows
\[
c^{\top}x-b^{\top}y=c^{\top}x-(\ma x)^{\top}y=c^{\top}x-x^{\top}(\ma^{\top}y)=x^{\top}s.
\]
Since $x,s\geq0$, this shows that 
\[
c^{\top}x\geq b^{\top}y.
\]
Optimizing over all $x$ and $y$, we obtain weak duality
\[
\min_{\ma\ox=b,\ox\in\R^{n}_{\geq0}}c^{\top}\ox\geq\max_{\ma^{\top}y+s=c,s\in\R^{n}_{\geq0}}b^{\top}y.
\]
Using this, we have
\[
c^{\top}x=b^{\top}y+x^{\top}s\leq\max_{\ma^{\top}\oy+\os=c,\os\in\R^{n}_{\geq0}}b^{\top}\oy+x^{\top}s\leq\min_{\ma\ox=b,\ox\in\R^{n}_{\geq0}}c^{\top}\ox+x^{\top}s.
\]
The other direction is similar.
\end{proof}
The main implication of Lemma \ref{lem:duality_gap} is that any feasible
$(x,s)$ with small $x^{\top}s$ is a nearly optimal solution of the
linear program. This leads us to primal-dual algorithms in which we
start with a feasible primal-dual solution pair $(x,s)$ and iteratively
update the solution to decrease the duality gap $x^{\top}s$. 

\section{Interior Point Method}

In this section, we discuss the classical short-step interior point
method. For two vectors $a,b$, we use $ab$ to denote the vector
with components $(ab)_{i}=a_{i}b_{i}$ and $a/b$ to denote the vector
with components $a_{i}/b_{i}$. For a scalar $t\in\R$, we let $t\mathbf{1}$
denote the vector with all coordinates equal to $t$. 
\begin{defn}[Central Path]
\label{def:central_path}We define the central path $(x_{t},s_{t})\in\PriR^{\circ}\times\DualR^{\circ}$
by $x_{t}s_{t}=t\mathbf{1}$. We say $x_{t}$ is on the central path
of (\ref{eq:tut_primal_LP}) at $t$.
\end{defn}

The algorithm maintains a pair $(x,s)\in\PriR^{\circ}\times\DualR^{\circ}$
and a scalar $t>0$ satisfying the invariant 
\[
\|\frac{xs}{t}-\mathbf{1}\|_{2}\leq\frac{1}{4}.
\]
Note that the deviation from the central path is measured in $\ell_{2}$
norm. In each step, the algorithm decreases $t$ by a factor of $1-\Omega(n^{-1/2})$
while maintaining the invariant.

\subsection{Basic Property of a Step}

To see why there is a pair $(x,s)\in\PriR^{\circ}\times\DualR^{\circ}$
satisfying the invariant, we prove the following generalization.
\begin{lem}[Quadrant Representation of Primal-Dual]
\label{lem:primal_dual_representation}Suppose $\PriR^{\circ}$ is
non-empty and $\PriR$ is bounded. For any positive vector $\mu\in\R^{n}_{>0}$,
there is a unique pair $(x_{\mu},s_{\mu})\in\PriR^{\circ}\times\DualR^{\circ}$
such that $x_{\mu}s_{\mu}=\mu$. Furthermore, $x_{\mu}=\arg\min_{x\in\PriR}f_{\mu}(x)$
where
\[
f_{\mu}(x)=c^{\top}x-\sum^{n}_{i=1}\mu_{i}\ln x_{i}.
\]
\end{lem}

\begin{proof}
Fix $\mu\in\R^{n}_{>0}$. We define $x_{\mu}=\arg\min_{x\in\PriR}f_{\mu}(x)$
and prove that $(x_{\mu},s_{\mu})\in\PriR^{\circ}\times\DualR^{\circ}$
with $x_{\mu}s_{\mu}=\mu$ for some $s_{\mu}$. Since $\PriR^{\circ}$
is non-empty, $\PriR$ is bounded, and $f_{\mu}$ is strictly convex,
such an $x_{\mu}$ exists. Furthermore, since $f_{\mu}(x)\rightarrow+\infty$
as $x_{i}\rightarrow0$ for any $i$, we have that $x_{\mu}\in\PriR^{\circ}$.

By the KKT optimality condition for $f_{\mu}$, there is a vector
$y$ such that
\[
\nabla f_{\mu}(x_{\mu})=c-\frac{\mu}{x_{\mu}}=\ma^{\top}y.
\]
Define $s_{\mu}=\frac{\mu}{x_{\mu}}$. Then one can check that $s_{\mu}\in\DualR^{\circ}$
and $x_{\mu}s_{\mu}=\mu$.

For the uniqueness, if $(x,s)\in\PriR^{\circ}\times\DualR^{\circ}$
and $xs=\mu$, then $x$ satisfies the optimality condition for $f_{\mu}$.
Since $f_{\mu}$ is strictly convex, such $x$ must be unique.
\end{proof}
Lemma \ref{lem:primal_dual_representation} shows that any point in
$\PriR^{\circ}\times\DualR^{\circ}$ is uniquely represented by a
positive vector $\mu$. Interior point methods move $\mu$ uniformly
to $0$ while maintaining the corresponding $x_{\mu}$. Now we discuss
how to find $(x_{\mu},s_{\mu})$ from a nearby interior feasible point
$(x,s)$; namely, how to move $(x,s)$ to $(x+\delta_{x},s+\delta_{s})$
so that the result satisfies the system
\begin{align*}
(x+\delta_{x})(s+\delta_{s}) & =\mu,\\
\ma(x+\delta_{x}) & =b,\\
\ma^{\top}(y+\delta_{y})+(s+\delta_{s}) & =c,\\
(x+\delta_{x},s+\delta_{s}) & \in\R^{2n}_{>0}.
\end{align*}
Although the system above involves $y$, our approximate solution
does not need to know $y$. By ignoring the inequality constraint
and the second-order term $\delta_{x}\delta_{s}$ in the first equation
above, and using $\ma x=b$ and $\ma^{\top}y+s=c$ we can simplify
the system:
\begin{align}
xs+\ms\delta_{x}+\mx\delta_{s} & =\mu,\label{eq:LP_newton_step}\\
\ma\delta_{x} & =0,\nonumber \\
\ma^{\top}\delta_{y}+\delta_{s} & =0,\nonumber 
\end{align}
where $\mx$ and $\ms$ are the diagonal matrices with diagonals $x$
and $s$, respectively. In the following lemma, we show how to write
the step above using a projection matrix ($\mproj^{2}=\mproj$).
\begin{lem}
\label{lem:Newton_step}Suppose that $\ma$ has full row rank and
$(x,s)\in\PriR^{\circ}\times\DualR^{\circ}$. Then, the unique solution
to the linear system (\ref{eq:LP_newton_step}) is given by
\begin{align*}
\mx^{-1}\delta_{x} & =(\mi-\mproj)(\delta_{\mu}/(xs)),\\
\ms^{-1}\delta_{s} & =\mproj(\delta_{\mu}/(xs))
\end{align*}
where $\delta_{\mu}=\mu-xs$, the division by $xs$ is coordinatewise,
and $\mproj=\ms^{-1}\ma^{\top}(\ma\ms^{-1}\mx\ma^{\top})^{-1}\ma\mx$.
\end{lem}

\begin{proof}
Note that the step satisfies $\ms\delta_{x}+\mx\delta_{s}=\delta_{\mu}$.
Multiplying both sides by $\ma\ms^{-1}$ and using $\ma\delta_{x}=0$,
we have
\[
\ma\ms^{-1}\mx\delta_{s}=\ma\ms^{-1}\delta_{\mu}.
\]
Now we use that $\ma^{\top}\delta_{y}+\delta_{s}=0$ and get
\[
\ma\ms^{-1}\mx\ma^{\top}\delta_{y}=-\ma\ms^{-1}\delta_{\mu}.
\]
Since $\ma\in\R^{d\times n}$ has full row rank and $\ms^{-1}\mx$
is invertible, we have that $\ma\ms^{-1}\mx\ma^{\top}$ is invertible.
Hence, 
\[
\delta_{y}=-(\ma\ms^{-1}\mx\ma^{\top})^{-1}\ma\ms^{-1}\delta_{\mu}\mbox{\ensuremath{\qquad} and \ensuremath{\qquad}}\delta_{s}=\ma^{\top}(\ma\ms^{-1}\mx\ma^{\top})^{-1}\ma\ms^{-1}\delta_{\mu}.
\]
Putting this into $\ms\delta_{x}+\mx\delta_{s}=\delta_{\mu}$, we
have
\[
\delta_{x}=\ms^{-1}\delta_{\mu}-\ms^{-1}\mx\ma^{\top}(\ma\ms^{-1}\mx\ma^{\top})^{-1}\ma\ms^{-1}\delta_{\mu}.
\]
The result follows from the definition of $\mproj$.
\end{proof}

\subsection{Lower Bounding Step Size}

The efficiency of interior point methods depends on how large a step
we can take while staying within the domain. In the rest of this subsection,
we write $\mu=xs$ for the current complementarity vector. We first
study the step operators $(\mi-\mproj)$ and $\mproj$. The following
lemma implies that $\mproj$ is a nearly orthogonal projection matrix
when $\mu$ is close to a multiple of the all-ones vector. Hence,
the relative changes of $\mx^{-1}\delta_{x}$ and $\ms^{-1}\delta_{s}$
are essentially the orthogonal decomposition of the relative step
$\delta_{\mu}/\mu$. For a vector $u$ and a positive vector $v$,
we define $\norm u_{v}\defeq\sqrt{u^{\top}\Diag(v)u}$.
\begin{lem}
\label{lem:Projection_property}Under the assumption in Lemma \ref{lem:Newton_step},
$\mproj$ is a projection matrix such that $\|\mproj v\|_{\mu}\leq\|v\|_{\mu}$
for any $v\in\Rn$. Similarly, we have that $\|(\mi-\mproj)v\|_{\mu}\leq\|v\|_{\mu}$.
\end{lem}

\begin{proof}
$\mproj$ is a projection because $\mproj^{2}=\mproj$. Define the
orthogonal projection matrix
\[
\mproj_{\mathrm{orth}}=\ms^{-1/2}\mx^{1/2}\ma^{\top}(\ma\ms^{-1}\mx\ma^{\top})^{-1}\ma\mx^{1/2}\ms^{-1/2},
\]
then we have
\begin{align*}
\|\mproj v\|^{2}_{\mu} & =v^{\top}\mx\ma^{\top}(\ma\ms^{-1}\mx\ma^{\top})^{-1}\ma\ms^{-1}\mx\ms\ms^{-1}\ma^{\top}(\ma\ms^{-1}\mx\ma^{\top})^{-1}\ma\mx v\\
 & =v^{\top}\ms^{1/2}\mx^{1/2}\mproj_{\mathrm{orth}}\ms^{1/2}\mx^{1/2}v\\
 & \leq v^{\top}\ms^{1/2}\mx^{1/2}\ms^{1/2}\mx^{1/2}v=\|v\|^{2}_{\mu}.
\end{align*}
The calculation for $\|(\mi-\mproj)v\|_{\mu}$ is similar.
\end{proof}
Next we give a lower bound on the largest feasible step size.
\begin{lem}
\label{lem:step_size}We have that $\|\mx^{-1}\delta_{x}\|^{2}_{\infty}\leq\frac{1}{\min_{i}\mu_{i}}\|\delta_{\mu}/\mu\|^{2}_{\mu}$
and $\|\ms^{-1}\delta_{s}\|^{2}_{\infty}\leq\frac{1}{\min_{i}\mu_{i}}\|\delta_{\mu}/\mu\|^{2}_{\mu}$.
In particular, if $\|\delta_{\mu}/\mu\|^{2}_{\mu}<\min_{i}\mu_{i}$,
we have $(x+\delta_{x},s+\delta_{s})\in\PriR^{\circ}\times\DualR^{\circ}$.
\end{lem}

\begin{proof}
For $\|\mx^{-1}\delta_{x}\|_{\infty}$, we have $\min_{i}\mu_{i}\|\mx^{-1}\delta_{x}\|^{2}_{\infty}\leq\|\mx^{-1}\delta_{x}\|^{2}_{\mu}$
and hence
\[
\|\mx^{-1}\delta_{x}\|^{2}_{\infty}\leq\frac{1}{\min_{i}\mu_{i}}\|\mx^{-1}\delta_{x}\|^{2}_{\mu}=\frac{1}{\min_{i}\mu_{i}}\|(\mi-\mproj)(\delta_{\mu}/\mu)\|^{2}_{\mu}\leq\frac{1}{\min_{i}\mu_{i}}\|\delta_{\mu}/\mu\|^{2}_{\mu}.
\]
The proof for $\|\ms^{-1}\delta_{s}\|_{\infty}$ is similar.

Hence, if $\|\delta_{\mu}/\mu\|^{2}_{\mu}<\min_{i}\mu_{i}$, we have
that $\|\mx^{-1}\delta_{x}\|_{\infty}<1$ and $\|\ms^{-1}\delta_{s}\|_{\infty}<1$,
i.e., $|\delta_{x,i}|<|x_{i}|$ and $|\delta_{s,i}|<|s_{i}|$ for
all $i$. Therefore, $x+\delta_{x}>0$ and $s+\delta_{s}>0$ are feasible.
\end{proof}
To decrease $\mu$ uniformly, we set $\delta_{\mu}=-h\mu$ for some
step size $h$. To ensure feasibility, we need $\|\delta_{\mu}/\mu\|^{2}_{\mu}\leq\min_{i}\mu_{i}$
and this gives the maximum step size
\begin{equation}
h=\sqrt{\frac{\min_{i}\mu_{i}}{\sum_{i}\mu_{i}}}.\label{eq:max_h_formula}
\end{equation}
Note that the above quantity is maximized at $h=n^{-1/2}$ when $\mu$
has all equal coordinates.

\subsection{Staying within a small $\ell_{2}$ distance}

The step size (\ref{eq:max_h_formula}) is maximized when $\mu$ is
a constant vector. A natural approach is therefore to keep $\mu$
as a vector close in $\ell_{2}$ norm to a multiple of the all-ones
vector. This motivates the following algorithm:

\begin{algorithm2e}[H]

\caption{$\mathtt{L2Step}(\ma,x,s,t_{\mathrm{start}},t_{\mathrm{end}})$}

\label{alg:L2Step}

\textbf{Define} $\mproj_{x,s}=\ms^{-1}\ma^{\top}(\ma\ms^{-1}\mx\ma^{\top})^{-1}\ma\mx$.

\textbf{Invariant:} $(x,s)\in\PriR^{\circ}\times\DualR^{\circ}$ and
$\|xs-t\mathbf{1}\|_{2}\leq\frac{t}{4}$. 

Let $t=t_{\mathrm{start}}$ and $h=1/(16\sqrt{n})$, where $n$ is
the number of columns in $\ma$.

\Repeat{$t=t_{\mathrm{end}}$}{

Let $t'=\max(t/(1+h),t_{\mathrm{end}})$.

Let $\mu=xs$ and $\delta_{\mu}=t'\mathbf{1}-\mu$.

Let $\delta_{x}=\mx(\mi-\mproj_{x,s})(\delta_{\mu}/\mu)$ and $\delta_{s}=\ms\mproj_{x,s}(\delta_{\mu}/\mu)$.

Set $x\leftarrow x+\delta_{x}$, $s\leftarrow s+\delta_{s}$ and $t\leftarrow t'$.

}

\textbf{Return} $(x,s)$.

\end{algorithm2e}

Note that the algorithm requires some initial point $(x,s)$ close
to the central path and we will show how to get this in the next section
(by changing the linear program temporarily).

First, we show that the invariant is maintained in each step. The
conclusion that the distance is less than $t/6$ is needed in Appendix
\ref{subsec:find_IPM_interior}, where we call $\mathtt{L2Step}$
on a modified LP and then prove that the result is close to the central
path for the original LP.
\begin{lem}
\label{lem:L2Step}Suppose that the input satisfies $(x,s)\in\PriR^{\circ}\times\DualR^{\circ}$
and $\|xs-t_{\mathrm{start}}\mathbf{1}\|_{2}\leq\frac{t_{\mathrm{start}}}{4}$.
Then, the algorithm $\mathtt{L2Step}$ maintains $(x,s)\in\PriR^{\circ}\times\DualR^{\circ}$
and $t>0$ such that $\|xs-t\mathbf{1}\|_{2}\leq\frac{t}{6}$.
\end{lem}

\begin{proof}
We prove by induction that $\|xs-t\mathbf{1}\|_{2}\leq\frac{t}{6}$
after each step. Note that the input satisfies $\|xs-t\mathbf{1}\|_{2}\leq\frac{t}{4}$.

Let $x'=x+\delta_{x}$, $s'=s+\delta_{s}$ and let $t'$ be as defined
in the algorithm. Note that
\[
x's'-t'\mathbf{1}=(x+\delta_{x})(s+\delta_{s})-t'\mathbf{1}=\mu+\ms\delta_{x}+\mx\delta_{s}+\delta_{x}\delta_{s}-t'\mathbf{1}.
\]
Lemma \ref{lem:Newton_step} shows that $\ms\delta_{x}+\mx\delta_{s}=t'\mathbf{1}-\mu$.
Hence, we have 
\[
x's'-t'\mathbf{1}=\delta_{x}\delta_{s}=\mx^{-1}\delta_{x}\cdot\ms^{-1}\delta_{s}\cdot\mu.
\]
Using this, we have
\[
\|x's'-t'\mathbf{1}\|_{2}\leq\|\mu^{1/2}\mx^{-1}\delta_{x}\|_{2}\|\mu^{1/2}\ms^{-1}\delta_{s}\|_{2}=\|\mx^{-1}\delta_{x}\|_{\mu}\|\ms^{-1}\delta_{s}\|_{\mu}\leq\|\delta_{\mu}/\mu\|^{2}_{\mu}.
\]
Here we used Lemma \ref{lem:Projection_property} at the end.

Using $t'\mathbf{1}-\mu=\frac{t'}{t}(t\mathbf{1}-\mu)+(\frac{t'}{t}-1)\mu$,
we have
\[
\|\delta_{\mu}/\mu\|_{\mu}=\|\frac{t'}{t}\frac{t\mathbf{1}-\mu}{\mu}+(\frac{t'}{t}-1)\|_{\mu}\leq\frac{t'}{t}\|xs-t\mathbf{1}\|_{\mu^{-1}}+\|\frac{t'}{t}-1\|_{\mu}.
\]
Since $\|\mu-t\mathbf{1}\|_{2}\leq\frac{t}{4}$, we have $\min_{i}\mu_{i}\geq\frac{3t}{4}$
and $\max_{i}\mu_{i}\leq\frac{5}{4}t$. Using $|\frac{t'}{t}-1|\leq h=\frac{1}{16\sqrt{n}}$,
we have
\[
\|\delta_{\mu}/\mu\|_{\mu}\leq\frac{t'}{t}\sqrt{\frac{4}{3t}}\|xs-t\mathbf{1}\|_{2}+h\sqrt{\frac{5}{4}tn}\leq\sqrt{\frac{t}{12}}+h\sqrt{\frac{5}{4}tn}\leq0.36\sqrt{t}.
\]
Hence, we have $\|x's'-t'\mathbf{1}\|_{2}\leq\|\delta_{\mu}/\mu\|^{2}_{\mu}\leq0.13t\leq t'/6$.
Furthermore, $\|\delta_{\mu}/\mu\|^{2}_{\mu}<\min_{i}\mu_{i}$, which
implies $(x',s')$ is feasible (Lemma \ref{lem:step_size}).
\end{proof}
Note that the lemma above only concludes that the output is close
to the central path. To upper bound the error, we can apply Lemma
\ref{lem:duality_gap} which shows the duality gap is equal to $x^{\top}s$.

\subsection{Finding a feasible interior point by lifting\label{subsec:find_IPM_interior}}

Here we discuss how to get a feasible interior point close to the
central path by modifying the linear program. The runtime of an interior
point method depends on how degenerate the linear program is.
\begin{defn}
\label{def:LP_para}We define the following parameters for the linear
program $\min_{\ma x=b,x\geq0}c^{\top}x$:
\begin{enumerate}
\item Inner radius $r$: There exists an $x\in\PriR$ such that $x_{i}\geq r$
for all $i\in[n]$.
\item Outer radius $R$: For any $x\geq0$ with $\ma x=b$, we have that
$\|x\|_{2}\leq R$.
\item Lipschitz constant $L$: $\|c\|_{2}\leq L$.
\end{enumerate}
\end{defn}

Since $\mathtt{L2Step}$ requires a feasible point near the central
path, we modify the linear program to make it happen. To satisfy the
constraint $\ma x=b$, we start the algorithm by taking a least-squares
solution of the constraint $\ma x=b$. Since this solution can be
negative, we write the variable $x=x^{+}-x^{-}$ with both $x^{+},x^{-}\geq0$.
We put a large cost vector on $x^{-}$ to ensure the solution is roughly
the same. The crux of the proof is that if we optimize this new program
well enough, we will have $x^{+}-x^{-}>0$ and hence $x^{+}-x^{-}$
gives a good starting point for the original program. For technical
reasons, we need to put an extra constraint $1^{\top}x^{+}\leq\Lambda$
for some $\Lambda$ to ensure the problem is bounded.
\begin{thm}
\label{thm:modified_LP_construction}Given a linear program $\min_{\ma x=b,x\in\R^{n}_{\geq0}}c^{\top}x$
with inner radius $r$, outer radius $R$ and Lipschitz constant $L$
and any $0\leq\eta\leq\frac{1}{3}$. Consider the following modified
linear program:
\[
\min_{(x^{+},x^{-},x^{\theta})\in\PriR_{\eta}}\left\langle c^{+},x^{+}\right\rangle +\left\langle c^{-},x^{-}\right\rangle 
\]
with
\[
\PriR_{\eta}=\left\{ (x^{+},x^{-},x^{\theta})\in\R^{2n+1}_{\geq0}:\ \ma(x^{+}-x^{-})=b,\sum^{n}_{i=1}x^{+}_{i}+x^{\theta}=\widetilde{b}\right\} 
\]
where $\overline{R}=6R/\eta$, $\overline{t}=6000n^{2}R/(\eta^{2}r)\cdot LR$,
$\ox^{+}=\overline{t}/(c+\overline{t}/\overline{R})$, $\ox^{-}=\ox^{+}-\ma^{\top}(\ma\ma^{\top})^{-1}b$,
$c^{+}=c$, $c^{-}=\overline{t}/\ox^{-}$ and $\widetilde{b}=\sum^{n}_{i=1}\ox^{+}_{i}+\overline{R}$.
Then, we have that
\begin{enumerate}
\item $(\ox^{+},\ox^{-},\overline{R})$ is on the central path of the modified
program at $\overline{t}$.
\item For any primal $x\defeq(x^{+},x^{-},x^{\theta})\in\PriR_{\eta}$ and
dual $s\defeq(s^{+},s^{-},s^{\theta})$ such that $\frac{5}{6}LR\leq x_{i}s_{i}\leq\frac{7}{6}LR$
for all $i$, we have that
\[
(x^{+}-x^{-},s^{+}-s^{\theta})\in\PriR\times\DualR
\]
and that $x^{-}_{i}\leq\eta x^{+}_{i}$ and $s^{\theta}\leq\eta s^{+}_{i}$
for all $i$.
\end{enumerate}
\end{thm}

We continue from the discussion in Section \ref{subsec:SolveLP_apx}.

We first show that $(\ox^{+},\ox^{-},\overline{R})$ defined in Theorem
\ref{thm:modified_LP_construction} is indeed on the central path
of the modified linear program. We define the dual of the modified
linear program as follows
\[
\mathcal{D}_{\eta}=\{(s^{+},s^{-},s^{\theta})\in\R^{2n+1}_{\geq0}:\ma^{\top}y+\lambda\mathbf{1}+s^{+}=c^{+},-\ma^{\top}y+s^{-}=c^{-},\lambda+s^{\theta}=0\text{ for some }y\in\R^{d}\text{ and }\lambda\in\R\}.
\]

\begin{lem}
\label{lem:ipm_interior_modified}The modified linear program in Theorem
\ref{thm:modified_LP_construction} has an explicit point on the central
path,$(\ox^{+},\ox^{-},\overline{R})$ at $\overline{t}$.
\end{lem}

\begin{proof}
Recall that we say $(x^{+},x^{-},x^{\theta})$ is on the central path
at $t$ if $x^{+},x^{-},x^{\theta}$ are positive, $x^{+}_{i}s^{+}_{i}=t,x^{-}_{i}s^{-}_{i}=t,x^{\theta}s^{\theta}=t$
and the following equations hold
\begin{align}
\ma x^{+}-\ma x^{-} & =b,\nonumber \\
\sum^{n}_{i=1}x^{+}_{i}+x^{\theta} & =\widetilde{b},\nonumber \\
\ma^{\top}y+\lambda+s^{+} & =c^{+},\label{eq:KKT_modified_LP}\\
-\ma^{\top}y+s^{-} & =c^{-},\nonumber \\
\lambda+s^{\theta} & =0,\nonumber 
\end{align}
for some $s^{+},s^{-}\in\R^{n}_{>0}$, $s^{\theta}>0$, $y\in\R^{d}$
and $\lambda\in\R$.

Now, we verify the solution $\ox^{+}=\overline{t}/(c+\overline{t}/\overline{R})$,
$\ox^{-}=\ox^{+}-\ma^{\top}(\ma\ma^{\top})^{-1}b$, $\ox^{\theta}=\overline{R}$,
$y=0$, $\os^{+}=\overline{t}/\ox^{+}$, $\os^{-}=\overline{t}/\ox^{-}$,
$\os^{\theta}=\overline{t}/\ox^{\theta}$, $\lambda=-\os^{\theta}$.
Using $\ma\ma^{\top}(\ma\ma^{\top})^{-1}b=b$, one can check that
it satisfies all the equality constraints above. 

For the positivity constraints, using $\|c\|_{\infty}\leq L$ and
$\overline{t}\geq4L\overline{R}$, we have
\begin{equation}
\frac{4}{5}\overline{R}\leq\frac{\overline{t}}{L+\overline{t}/\overline{R}}\leq\ox^{+}_{i}\leq\frac{\overline{t}}{-L+\overline{t}/\overline{R}}\leq\frac{4}{3}\overline{R}\label{eq:x_pos_bound}
\end{equation}
and hence $\ox^{+}>0$ and so is $\os^{+}$. Note that $x_{\circ}\defeq\ma^{\top}(\ma\ma^{\top})^{-1}b$
is the least-squares solution to $\ma x=b$, so, we have 
\begin{equation}
\|x_{\circ}\|_{2}\leq\|x^{*}\|_{2}\leq R<\frac{1}{4}\overline{R}.\label{eq:x0_upper_bound}
\end{equation}
Together with $\ox^{+}_{i}\geq\frac{4}{5}\overline{R}$, we have
\[
\ox^{-}_{i}\geq\frac{4}{5}\overline{R}-\frac{1}{4}\overline{R}>0
\]
for all $i$. Hence, $\ox^{-}$ and $\os^{-}$ are positive. Finally,
$\ox^{\theta}$ and $\os^{\theta}$ are positive. This proves that
$(\ox^{+},\ox^{-},\overline{R})$ is on the central path at $\overline{t}$.
\end{proof}
Next, we show that the near-central-path point $(x,s)$ at $t=LR$
is far from the boundary $x^{+}=0$ and is close to the boundary $x^{-}=0$.
The proofs of both statements use the same idea: the optimality condition
of $x$. Throughout the rest of the proof, we are given $(x,s)\in\PriR_{\eta}\times\mathcal{D}_{\eta}$
such that $\mu=xs$ satisfies
\[
\frac{5}{6}LR\leq\mu\leq\frac{7}{6}LR.
\]
We write $\mu$ into its three parts $(\mu^{+},\mu^{-},\mu^{\theta})$.
By Lemma \ref{lem:primal_dual_representation}, we have that $x\defeq(x^{+},x^{-},x^{\theta})$
minimizes the function
\[
f_{\mu}(x^{+},x^{-},x^{\theta})\defeq\left\langle c^{+},x^{+}\right\rangle +\left\langle c^{-},x^{-}\right\rangle -\sum^{n}_{i=1}\mu^{+}_{i}\log x^{+}_{i}-\sum^{n}_{i=1}\mu^{-}_{i}\log x^{-}_{i}-\mu^{\theta}\log x^{\theta}
\]
over the domain $\PriR_{\eta}$. The gradient of $f_{\mu}$ is somewhat
complicated. We avoid it by considering the directional derivative
at $x$ along a feasible direction toward a lifted version of a point
$v$ such that $\ma v=b$ and $v\geq r$. Since our domain is in $\PriR_{\eta}\subset\R^{2n+1}$,
we lift $v$ to the higher-dimensional point defined by
\begin{align*}
v^{-} & =\min(x^{-},\frac{4L\overline{R}}{\overline{t}}\cdot R),\\
v^{+} & =v+v^{-},\\
v^{\theta} & =\widetilde{b}-\sum^{n}_{i=1}v^{+}_{i}.
\end{align*}
First, we need to get some basic bounds on $\widetilde{b}$ and $c^{-}$.
\begin{lem}
\label{lem:ipm_propert_modified}We have that $\frac{4}{5}n\overline{R}\leq\widetilde{b}\leq3n\overline{R}$
and $c^{-}_{i}\geq3\overline{t}/(7\overline{R})$ for all $i$.
\end{lem}

\begin{proof}
By (\ref{eq:x_pos_bound}), we have $\frac{4}{5}\overline{R}\leq\ox^{+}_{i}\leq\frac{4}{3}\overline{R}$.
By the definition of $\widetilde{b}$, we have
\[
\widetilde{b}=\sum^{n}_{i=1}\ox^{+}_{i}+\overline{R}\leq\frac{4}{3}n\overline{R}+\overline{R}\leq3n\overline{R}.
\]
Similarly, we have $\widetilde{b}=\sum_{i}\ox^{+}_{i}+\overline{R}\geq\frac{4}{5}n\overline{R}$. 

For the bound on $c^{-}_{i}$, recall that $c^{-}_{i}=\overline{t}/\ox^{-}_{i}$
with $\ox^{-}=\ox^{+}-x_{\circ}$ and $x_{\circ}=\ma^{\top}(\ma\ma^{\top})^{-1}b$.
Using (\ref{eq:x0_upper_bound}), we have 
\[
\ox^{-}_{i}\leq\frac{4}{3}\overline{R}+R.
\]
Therefore, $c^{-}\geq3\overline{t}/(7\overline{R})$.
\end{proof}
The following lemma shows that $(v^{+},v^{-},v^{\theta})\in\PriR_{\eta}$.
\begin{lem}
\label{lem:IPM_exact_inner_pt}We have that $(v^{+},v^{-},v^{\theta})\in\PriR_{\eta}$.
Furthermore, we have $v^{\theta}\geq\frac{1}{2}n\overline{R}$.
\end{lem}

\begin{proof}
Note that $(v^{+},v^{-},v^{\theta})$ satisfies the linear constraints
of $\PriR_{\eta}$ by construction. It suffices to prove that the
vector is positive. Since $x^{-}>0$, we have $v^{-}>0$. Since $v\geq r$,
we also have $v^{+}>0$. For $v^{\theta}$, we use $\widetilde{b}\geq\frac{4}{5}n\overline{R}$
(Lemma \ref{lem:ipm_propert_modified}), $v_{i}\leq R$ and $v^{-}\leq\frac{4L\overline{R}}{t}\cdot R\leq R$
to get
\begin{equation}
v^{+}_{i}=v_{i}+v^{-}_{i}\leq2R\label{eq:v_plus_upper}
\end{equation}
and hence
\begin{align*}
v^{\theta}=\widetilde{b}-\sum^{n}_{i=1}v^{+}_{i} & \geq\frac{4}{5}n\overline{R}-\sum^{n}_{i=1}(v_{i}+v^{-}_{i})\geq\frac{1}{2}n\overline{R}.
\end{align*}
\end{proof}
Next, we define the path $p(\tau)=(1-\tau)(x^{+},x^{-},x^{\theta})+\tau(v^{+},v^{-},v^{\theta})$.
Since $p(0)$ minimizes $f_{\mu}$, we have that $\frac{d}{d\tau}f_{\mu}(p(\tau))|_{\tau=0}\geq0$.
In particular, we have
\begin{align}
0 & \leq\frac{d}{d\tau}f_{\mu}(p(\tau))|_{\tau=0}\nonumber \\
 & =\left\langle c,v^{+}-x^{+}\right\rangle +\left\langle c^{-},v^{-}-x^{-}\right\rangle -\sum^{n}_{i=1}\frac{\mu^{+}_{i}}{x^{+}_{i}}(v^{+}-x^{+})_{i}-\sum^{n}_{i=1}\frac{\mu^{-}_{i}}{x^{-}_{i}}(v^{-}-x^{-})_{i}-\frac{\mu^{\theta}}{x^{\theta}}(v^{\theta}-x^{\theta})\nonumber \\
 & =\frac{\mu^{\theta}}{x^{\theta}}(x^{\theta}-v^{\theta})+\sum^{n}_{i=1}(c_{i}-\frac{\mu^{+}_{i}}{x^{+}_{i}})(v^{+}-x^{+})_{i}+\sum^{n}_{i=1}(c^{-}_{i}-\frac{\mu^{-}_{i}}{x^{-}_{i}})(v^{-}-x^{-})_{i}.\label{eq:LP_exact_3_term}
\end{align}
We now bound the three terms. For the first term, we note that 
\begin{equation}
\frac{\mu^{\theta}}{x^{\theta}}(x^{\theta}-v^{\theta})\leq\mu^{\theta}\leq2LR.\label{eq:LP_exact_3_term_1}
\end{equation}
For the second term in (\ref{eq:LP_exact_3_term}), we have the following:
\begin{lem}
\label{lem:LP_exact_3_term_1}We have that $\sum^{n}_{i=1}(c_{i}-\frac{\mu^{+}_{i}}{x^{+}_{i}})(v^{+}-x^{+})_{i}\leq4nLR-\frac{LRr}{2\min_{i}x^{+}_{i}}$.
\end{lem}

\begin{proof}
Note that
\begin{align*}
\sum^{n}_{i=1}(c_{i}-\frac{\mu^{+}_{i}}{x^{+}_{i}})(v^{+}-x^{+})_{i} & =\sum^{n}_{i=1}(c_{i}v^{+}_{i}-\frac{\mu^{+}_{i}}{x^{+}_{i}}v^{+}_{i}-c_{i}x^{+}_{i}+\mu^{+}_{i})\\
 & \leq\sum^{n}_{i=1}c_{i}v^{+}_{i}+\sum^{n}_{i=1}\mu^{+}_{i}-\sum^{n}_{i=1}\frac{\mu^{+}_{i}}{x^{+}_{i}}v^{+}_{i}\\
 & \leq\|c\|_{\infty}\|v^{+}\|_{1}+2nLR-\frac{1}{2}\sum^{n}_{i=1}\frac{LRr}{x^{+}_{i}}
\end{align*}
where we used $\mu^{+}_{i}\in[\frac{LR}{2},2LR]$ and $v^{+}_{i}\geq v_{i}\geq r$
at the end. The result follows from $\|c\|_{\infty}\leq L$, $\|v^{+}\|_{1}\leq2nR$
(\ref{eq:v_plus_upper}).
\end{proof}
For the third term in (\ref{eq:LP_exact_3_term}), we have the following:
\begin{lem}
\label{lem:LP_exact_3_term_2}We have that $\sum^{n}_{i=1}(c^{-}_{i}-\frac{\mu^{-}_{i}}{x^{-}_{i}})(v^{-}-x^{-})_{i}\leq\frac{1}{2}LR-\frac{\overline{t}}{8\overline{R}}\max_{i}x^{-}_{i}$.
\end{lem}

\begin{proof}
Using $v^{-}=\min(x^{-},\frac{4L\overline{R}}{\overline{t}}\cdot R)$,
we have $v^{-}_{i}\leq x^{-}_{i}$. We can ignore the terms with $v^{-}_{i}=x^{-}_{i}$. 

For $v^{-}_{i}<x^{-}_{i}$, we have $x^{-}_{i}\geq\frac{4L\overline{R}}{\overline{t}}R$.
Using $c^{-}_{i}\geq3\overline{t}/(7\overline{R})$ (Lemma \ref{lem:ipm_propert_modified}),
we have
\[
c^{-}_{i}-\frac{\mu^{-}_{i}}{x^{-}_{i}}\geq c^{-}_{i}-\frac{\mu^{-}_{i}}{\frac{4L\overline{R}}{\overline{t}}R}\geq c^{-}_{i}-\frac{\frac{7}{6}LR}{\frac{4L\overline{R}}{\overline{t}}R}=c^{-}_{i}-\frac{7\overline{t}}{24\overline{R}}\geq\frac{\overline{t}}{8\overline{R}}.
\]
Hence, we have
\[
\sum^{n}_{i=1}(c^{-}_{i}-\frac{\mu^{-}_{i}}{x^{-}_{i}})(v^{-}-x^{-})_{i}\leq\frac{\overline{t}}{8\overline{R}}\sum^{n}_{i=1}(v^{-}-x^{-})_{i}\leq\frac{\overline{t}}{8\overline{R}}(\frac{4L\overline{R}}{\overline{t}}\cdot R-\max_{i}x^{-}_{i}).
\]
\end{proof}
Combining (\ref{eq:LP_exact_3_term}), (\ref{eq:LP_exact_3_term_1}),
Lemma \ref{lem:LP_exact_3_term_1} and Lemma \ref{lem:LP_exact_3_term_2},
we have
\begin{align*}
0 & \leq\frac{\mu^{\theta}}{x^{\theta}}(x^{\theta}-v^{\theta})+\sum^{n}_{i=1}(c_{i}-\frac{\mu^{+}_{i}}{x^{+}_{i}})(v^{+}-x^{+})_{i}+\sum^{n}_{i=1}(c^{-}_{i}-\frac{\mu^{-}_{i}}{x^{-}_{i}})(v^{-}-x^{-})_{i}\\
 & \leq2LR+4nLR-\frac{LRr}{2\min_{i}x^{+}_{i}}+\frac{1}{2}LR-\frac{\overline{t}}{8\overline{R}}\max_{i}x^{-}_{i}\\
 & \leq6nLR-\frac{LRr}{2\min_{i}x^{+}_{i}}-\frac{\overline{t}}{8\overline{R}}\max_{i}x^{-}_{i}.
\end{align*}
Hence, we show that $(x^{+},x^{-},x^{\theta})$ is close to the central
path at $t=LR$ implies that $\min_{i}x^{+}_{i}$ cannot be too small
and $\max_{i}x^{-}_{i}$ cannot be too large
\[
\frac{LRr}{2\min_{i}x^{+}_{i}}+\frac{\overline{t}}{8\overline{R}}\max_{i}x^{-}_{i}\leq6nLR.
\]
In particular, this shows the following:
\begin{lem}
\label{lem:LP_exact_distance_OPT}We have that $\min_{i}x^{+}_{i}\geq\frac{r}{12n}$
and $\max_{i}x^{-}_{i}\leq\frac{48nL\overline{R}}{\overline{t}}\cdot R$.
\end{lem}

Now, we are ready to prove Theorem \ref{thm:modified_LP_construction}.
\begin{proof}[Proof of Theorem \ref{thm:modified_LP_construction}]
First, we check $x\defeq x^{+}-x^{-}\in\PriR$. By the choice of
$\overline{R}$ and $\overline{t}$, Lemma \ref{lem:LP_exact_distance_OPT}
shows that 
\[
\frac{\max_{i}x^{-}_{i}}{\min_{i}x^{+}_{i}}\leq\frac{\frac{48nL\overline{R}}{\overline{t}}\cdot R}{\frac{r}{12n}}=576n^{2}\frac{R}{r}\cdot\frac{L\overline{R}}{\overline{t}}<\eta
\]
Hence, we have $x^{+}-x^{-}>0$ and that $\ma(x^{+}-x^{-})=b$. 

Next, we check $s\defeq s^{+}-s^{\theta}\in\DualR$. Since $x\in\mathcal{P}$,
we have $x_{i}\leq R$. Since $x^{-}_{i}\leq\frac{1}{2}x^{+}_{i}$,
we have $x^{+}_{i}\leq2x_{i}\leq2R$. Since $x^{+}_{i}s^{+}_{i}\geq\frac{5}{6}LR$,
we have $s^{+}_{i}\geq\frac{1}{3}L.$ On the other hand, we have $x^{\theta}=\widetilde{b}-\sum^{n}_{i=1}x^{+}_{i}\geq\widetilde{b}-2nR\geq\frac{1}{2}n\overline{R}$
(Lemma \ref{lem:ipm_propert_modified}). Hence, $s^{\theta}\leq\frac{\frac{7}{6}LR}{\frac{1}{2}n\overline{R}}\leq\frac{3LR}{n\overline{R}}.$
Combining both and the choice of $\overline{R}$, we have
\[
\frac{s^{\theta}}{\min_{i}s^{+}_{i}}\leq\frac{\frac{3LR}{n\overline{R}}}{L/3}=\frac{9R}{n\overline{R}}<\eta
\]
Hence, we have $s=s^{+}-s^{\theta}>0$ and that $\ma^{\top}y+s=\ma^{\top}y+s^{+}-s^{\theta}=\ma^{\top}y+s^{+}+\lambda=c$
(See (\ref{eq:KKT_modified_LP})).
\end{proof}
To ensure the reduction does not increase the complexity of solving
the linear systems, we note that the linear constraint in the modified
linear program is
\[
\oma=\left[\begin{array}{ccc}
\ma & -\ma & 0\\
1^{\top}_{n} & 0 & 1
\end{array}\right]
\]
For any diagonal matrices $\mw_{1},\mw_{2}$ with diagonals $w_{1},w_{2}$
and any scalar $\alpha$, we have
\[
\mh\defeq\oma\left[\begin{array}{ccc}
\mw_{1} & 0 & 0\\
0 & \mw_{2} & 0\\
0 & 0 & \alpha
\end{array}\right]\oma^{\top}=\left[\begin{array}{cc}
\ma(\mw_{1}+\mw_{2})\ma^{\top} & \ma w_{1}\\
(\ma w_{1})^{\top} & \|w_{1}\|_{1}+\alpha
\end{array}\right].
\]
Note that the second row and second column blocks have size $1$.
By the block inverse formula, $\mh^{-1}v$ is an explicit formula
involving $(\ma(\mw_{1}+\mw_{2})\ma^{\top})^{-1}v_{1:d}$ and $(\ma(\mw_{1}+\mw_{2})\ma^{\top})^{-1}\ma w_{1}$.
Hence, we can compute $\mh^{-1}v$ by solving two linear systems of
the form $\ma\mw\ma^{\top}$ with some extra linear work.

\subsection{Solving LP Approximately and Exactly\label{subsec:SolveLP_apx}}

Now we state our main algorithm.

\begin{algorithm2e}[H]

\caption{$\mathtt{SlowSolveLP}(\ma,b,c,\delta)$}

\label{alg:SlowSolveLP}

\textbf{Assumption: the linear program has inner radius $r$, outer
radius $R$ and Lipschitz constant $L$.}

Let $\eta=1/(48\sqrt{n})$, $\overline{R}=6R/\eta$, $\overline{t}=6000n^{2}R/(\eta^{2}r)\cdot LR$. 

\tcp{Define the modified program $\min_{\oma x=\overline{b}}\overline{c}^{\top}x$
by Theorem \ref{thm:modified_LP_construction} with parameters $\overline{R}$
and $\overline{t}$.}

Let $\oma=\left[\begin{array}{ccc}
\ma & -\ma & 0\\
1^{\top}_{n} & 0 & 1
\end{array}\right]$, $\overline{c}=(c,c^{-},0)$, $\overline{b}=(b,\widetilde{b})$ where
$c^{-}$ and $\widetilde{b}$ are defined in Theorem \ref{thm:modified_LP_construction}.

Let $(\overline{x},\overline{s})$ be the explicit central path at
$\overline{t}$ given by Theorem \ref{thm:modified_LP_construction}.

$(\overline{x},\overline{s})\leftarrow\mathtt{L2Step}(\oma,\overline{x},\overline{s},\overline{t},LR)$.

$(x,s)\leftarrow(\overline{x}^{+}-\overline{x}^{-},\overline{s}^{+}-\overline{s}^{\theta})$
where $\overline{x}=(\overline{x}^{+},\overline{x}^{-},\overline{x}^{\theta})$
and $\overline{s}=(\overline{s}^{+},\overline{s}^{-},\overline{s}^{\theta})$.

$(x_{\mathrm{end}},s_{\mathrm{end}})=\mathtt{L2Step}(\ma,x,s,LR,t_{\mathrm{end}})$
with $t_{\mathrm{end}}=\delta LR/(2n)$.

\textbf{Return} $x_{\mathrm{end}}$.

\end{algorithm2e}
\begin{thm}
\label{thm:SlowSolveLP}Consider a linear program $\min_{\ma x=b,x\geq0}c^{\top}x$
with $n$ variables and $d$ constraints. Assume the linear program
has inner radius $r$, outer radius $R$ and Lipschitz constant $L$
(see Definition \ref{def:LP_para}). Then, $\mathtt{SlowSolveLP}$
outputs $x$ such that
\begin{align*}
c^{\top}x & \leq\min_{\ma x=b,x\geq0}c^{\top}x+\delta LR,\\
\ma x & =b,\\
x & \geq0.
\end{align*}
The algorithm takes $O(\sqrt{n}\log(nR/(\delta r)))$ Newton steps
(defined in (\ref{eq:LP_newton_step})).

If we further assume that the solution $x^{*}=\arg\min_{\ma x=b,x\geq0}c^{\top}x$
is unique and that $c^{\top}x\geq c^{\top}x^{*}+\zeta LR$ for any
other vertex $x$ of $\{\ma x=b,x\geq0\}$ for some $\zeta>\delta\geq0$,
then we have that $\|x-x^{*}\|_{2}\leq\frac{2\delta R}{\zeta}$.
\end{thm}

\begin{proof}
By Theorem \ref{thm:modified_LP_construction}, the point $(\overline{x},\overline{s})$
is on the central path of the modified program at $\overline{t}$.
After the first call of $\mathtt{L2Step}$, Lemma \ref{lem:L2Step}
shows that $\mathtt{L2Step}$ returns $(\overline{x},\overline{s})$
such that $\|\overline{x}\overline{s}-t\|_{2}\leq\frac{t}{6}$ with
$t=LR$.

Theorem \ref{thm:modified_LP_construction} shows that $(x,s)=(\overline{x}^{+}-\overline{x}^{-},\overline{s}^{+}-\overline{s}^{\theta})\in\PriR\times\DualR$
and that $(1-\eta)\overline{x}^{+}_{i}\leq x_{i}\leq\overline{x}^{+}_{i}$
and $(1-\eta)\overline{s}^{+}_{i}\leq s_{i}\leq\overline{s}^{+}_{i}$
. Since $\eta=\frac{1}{48\sqrt{n}}$ and $\|\overline{x}^{+}\overline{s}^{+}-t\|_{2}\leq\frac{t}{6}$
with $t=LR$, we have that $\|xs-t\mathbf{1}\|_{2}\leq\frac{t}{4}$.
This verifies the condition for the second call of $\mathtt{L2Step}$.

After the second call of $\mathtt{L2Step}$, Lemma \ref{lem:L2Step}
shows that $\mathtt{L2Step}$ returns $(x_{\mathrm{end}},s_{\mathrm{end}})$
such that $\|x_{\mathrm{end}}s_{\mathrm{end}}-t_{\mathrm{end}}\|_{2}\leq\frac{t_{\mathrm{end}}}{6}$
with $t_{\mathrm{end}}=\delta LR/(2n)$. Hence, Lemma \ref{lem:duality_gap}
shows that 
\[
c^{\top}x_{\mathrm{end}}\leq\min_{\ma x=b,x\geq0}c^{\top}x+x^{\top}_{\mathrm{end}}s_{\mathrm{end}}\leq\min_{\ma x=b,x\geq0}c^{\top}x+2t_{\mathrm{end}}n\leq\min_{\ma x=b,x\geq0}c^{\top}x+\delta LR.
\]
For the runtime, note that $\mathtt{L2Step}$ decreases $t$ by a
factor of $1-\Omega(n^{-1/2})$ in each step. Hence, the first call
takes $O(\sqrt{n}\log(nR/r))$ Newton steps and the second call takes
$O(\sqrt{n}\log(n/\delta))$ Newton steps.

For the last conclusion, we assume $\delta\leq\zeta$ and, for $\alpha\geq0$,
let $\mathcal{P}(\alpha)=\mathcal{P}\cap\{c^{\top}x\leq c^{\top}x^{*}+\alpha LR\}$.
Note that $\mathcal{P}(\zeta)$ is a cone at $x^{*}$ (because there
is no vertex other than $x^{*}$ with value less than $c^{\top}x^{*}+\zeta LR$).
Hence, we have $\mathcal{P}(\delta)-x^{*}=\frac{\delta}{\zeta}(\mathcal{P}(\zeta)-x^{*})$.
Since $x\in\mathcal{P}(\delta)$, we have that 
\[
\|x-x^{*}\|_{2}\leq\frac{\delta}{\zeta}\text{diameter}(\mathcal{P}(\zeta)-x^{*})\leq\frac{2\delta R}{\zeta}.
\]
\end{proof}
If we know the solution of the linear program is integral or rational
with some bound on the number of bits, then getting a solution close
enough to $x^{*}$ allows us to round the solution to an integral/rational
solution. Therefore, the last conclusion of the theorem above gives
us an exact linear program algorithm assuming $\ma,b,c$ are integral
and bounded. The uniqueness assumption can be achieved by perturbing
the cost vector by a random vector.
\begin{xca}
Make the perturbation deterministic while preserving solutions.
\end{xca}


\subsection{Why $\sqrt{n}$?}

The central path is the solution to the following ODE
\begin{align*}
S_{t}\frac{d}{dt}x_{t}+X_{t}\frac{d}{dt}s_{t} & =1,\\
A\frac{d}{dt}x_{t} & =0,\\
A^{\top}\frac{d}{dt}y_{t}+\frac{d}{dt}s_{t} & =0.
\end{align*}
Solving this linear system, we have that $S_{t}\frac{dx_{t}}{dt}=(I-P_{t})1$
and $X_{t}\frac{ds_{t}}{dt}=P_{t}1$ where $P_{t}=X_{t}A^{\top}(AS^{-1}_{t}X_{t}A^{\top})^{-1}AS^{-1}_{t}$.
Using that $x_{t}s_{t}=t$, we have that
\[
P_{t}=X_{t}A^{\top}(AX^{2}_{t}A^{\top})^{-1}AX_{t}=S^{-1}_{t}A^{\top}(AS^{-2}_{t}A^{\top})^{-1}AS^{-1}_{t}
\]
and that $X^{-1}_{t}\frac{dx_{t}}{dt}=\frac{1}{t}(I-P_{t})1$ and
$S^{-1}_{t}\frac{ds_{t}}{dt}=\frac{1}{t}P_{t}1$. Equivalently, we
have
\[
\frac{d\ln x_{t}}{d\ln t}=(I-P_{t})1\text{ and }\frac{d\ln s_{t}}{d\ln t}=P_{t}1.
\]
Note that 
\[
\norm{P_{t}1}_{\infty}\leq\norm{P_{t}1}_{2}\leq\sqrt{n}.
\]
Hence, $x_{t}$ and $s_{t}$ can change by at most a constant factor
when we change $t$ by a $1\pm\frac{1}{\sqrt{n}}$ factor.
\begin{xca}
If we are given $x$ such that $\norm{\ln x-\ln x_{t}}_{\infty}=O(1)$,
then we can find $x_{t}$ by solving $\tilde{O}(1)$ linear systems.

\end{xca}


\section{Interior Point Method for Convex Programs}

The interior point method can be used to optimize any convex function.
For a more in-depth treatment, please see \cite{nesterov1998introductory}.

Recall that any convex optimization problem
\[
\min_{x}f(x)
\]
can be rewritten in the epigraph form as
\[
\min_{\{(x,t):\ f(x)\leq t\}}t.
\]
Hence, it suffices to study the problem $\min_{x\in K}c^{\top}x$.
Similar to the case of linear programs, we replace the hard constraint
$x\in K$ by a soft constraint as follows:
\[
\min_{x}\phi_{t}(x)\text{ where }\phi_{t}(x)=tc^{\top}x+\phi(x).
\]
where $\phi(x)$ is a convex function such that $\phi(x)\rightarrow+\infty$
as $x\rightarrow\partial K$. Note that we put the parameter $t$
in front of the cost $c^{\top}x$ instead of $\phi$ as in the linear-programming
formulation above; it is slightly more convenient here. We say $\phi$
is a \emph{barrier} for $K$. To be concrete, we can always keep in
mind $\phi(x)=-\sum^{n}_{i=1}\ln x_{i}$. As before, we define the
central path.
\begin{defn}
The central path $x_{t}=\arg\min_{x}\phi_{t}(x)$.
\end{defn}

The interior point method follows the following framework:
\begin{enumerate}
\item Find $x$ close to $x_{0}$.
\item While $t$ is not large enough,
\begin{enumerate}
\item Move $x$ closer to $x_{t}$
\item $t\rightarrow(1+h)\cdot t$.
\end{enumerate}
\end{enumerate}

\subsection{Self-concordance}

In this section, we give a general analysis for the Newton method.
In the next section, we will use this to show that interior point
method can be generalized to convex optimization. A key property of
the Newton method is that it is invariant under linear transformation.
In general, whenever a method uses $k^{\mathrm{th}}$ order information,
we need to assume the $k^{\mathrm{th}}$ derivative is continuous.
Otherwise, the $k^{\mathrm{th}}$ derivative is not useful for algorithmic
purposes. For the Newton method, it is convenient to assume that the
Hessian is Lipschitz. Since the method is invariant under linear transformation,
it only makes sense to impose an assumption that is invariant under
linear transformation.
\begin{defn}
Given a convex function $f:\R^{n}\rightarrow\R$, and any point $x\in\R^{n},$
define the norm $\norm{\cdot}_{x}$ as 
\[
\norm v^{2}_{x}=v^{\top}\nabla^{2}f(x)v.
\]
We call a function $f$ self-concordant if for any $h\in\Rn$ and
any $x$ in $\dom f$, we have
\[
|D^{3}f(x)[h,h,h]|\leq2\norm h^{3}_{x}
\]
where $D^{k}f(x)[h_{1},h_{2},\cdots,h_{k}]$ is the directional $k^{\mathrm{th}}$
derivative of $f$ along the directions $h_{1},h_{2},\cdots,h_{k}$.
\end{defn}

\begin{rem*}
The constant $2$ is chosen so that $-\ln(x)$ exactly satisfies the
assumption and it is not very important, in that by scaling $f$,
we can change any constant to any other constant.
\end{rem*}
\begin{xca}
Show that the following property is equivalent to self-concordance
as defined above: restricted on any straight line $g(t)=f(x+th)$,
we have $|g'''(t)|\leq2g''(t)^{3/2}$.
\end{xca}

\begin{xca}
Show that the functions $x^{\top}Ax$, $-\ln x$, $-\ln(1-\sum x^{2}_{i})$,
$-\ln\det X$ are self-concordant for $A\succeq0$.
\end{xca}

The self-concordance condition says that locally, the Hessian does
not change too fast, i.e., the change in the Hessian is bounded by
its magnitude (to the power $1.5$). We will skip the proof of the
lemma below.
\begin{lem}
Given a self-concordant function $f$, for any $h_{1},h_{2},h_{3}\in\Rn$,
we have 
\[
\left|D^{3}f(x)[h_{1},h_{2},h_{3}]\right|\leq2\norm{h_{1}}_{x}\norm{h_{2}}_{x}\norm{h_{3}}_{x}.
\]
\end{lem}

From the self-concordance condition, we have the following more directly
usable property.
\begin{lem}
\label{lem:self_con_global}For a self-concordant function $f$ and
any $x\in\dom f$ and any $\norm{y-x}_{x}<1$, we have that
\[
(1-\norm{y-x}_{x})^{2}\nabla^{2}f(x)\preceq\nabla^{2}f(y)\preceq\frac{1}{(1-\norm{y-x}_{x})^{2}}\nabla^{2}f(x).
\]
\end{lem}

\begin{proof}
Let $\alpha(t)=\left\langle \nabla^{2}f(x+t(y-x))u,u\right\rangle $.
Then, we have that 
\[
\alpha'(t)=D^{3}f(x+t(y-x))[y-x,u,u].
\]
By self-concordance, we have
\begin{equation}
\left|\alpha'(t)\right|\leq2\norm{y-x}_{x+t(y-x)}\norm u^{2}_{x+t(y-x)}.\label{eq:self_con_global_alpha}
\end{equation}

For $u=y-x$, we have $\left|\alpha'(t)\right|\leq2\alpha(t)^{\frac{3}{2}}$.
Hence, we have $\frac{d}{dt}\frac{1}{\sqrt{\alpha(t)}}\geq-1.$ Integrating
both sides with respect to $t$, we have
\[
\frac{1}{\sqrt{\alpha(t)}}\geq\frac{1}{\sqrt{\alpha(0)}}-t=\frac{1}{\norm{x-y}_{x}}-t.
\]
Rearranging it gives
\[
\norm{y-x}^{2}_{x+t(y-x)}\leq\frac{1}{(\frac{1}{\norm{x-y}_{x}}-t)^{2}}=\frac{\norm{x-y}^{2}_{x}}{(1-t\norm{x-y}_{x})^{2}}.
\]

Going back to general $u$, (\ref{eq:self_con_global_alpha}) gives
\[
\left|\alpha'(t)\right|\leq2\frac{\norm{x-y}_{x}}{1-t\norm{x-y}_{x}}\alpha(t).
\]
Rearranging,
\[
\left|\frac{d}{dt}\ln\alpha(t)\right|\leq2\frac{\norm{x-y}_{x}}{1-t\norm{x-y}_{x}}=-2\frac{d}{dt}\ln(1-t\norm{x-y}_{x})
\]
Integrating both from $t=0$ to $1$ gives the result.
\end{proof}

\subsection{Self-concordant barrier functions}

To analyze the algorithm above, we need to assume that $\phi$ is
well-behaved. We measure the quality of $\phi$ by $\|\nabla\phi\|_{\nabla^{2}\phi(x)^{+}}$,
where $A^{+}$ denotes the \emph{pseudo inverse} of $A$. If $A$
is invertible, then $A^{+}=A^{-1}$. One can think of this as the
Lipschitz constant of $\phi$ but measured in the local norm. 
\begin{defn}
We call $\phi$ a $\nu$-self-concordant barrier for $K$ if $\phi$
is self-concordant, $\phi(x)\rightarrow+\infty$ as $x\rightarrow\partial K$
and that $\|\nabla\phi(x)\|^{2}_{\nabla^{2}\phi(x)^{+}}\leq\nu$ for
all $x$.
\end{defn}

\begin{rem*}
We use the pseudo-inverse because we only care about the inverse in
the direction of the gradient: the above definition only makes sense
if, for any $x\in K$, there is a vector $v_{x}$ such that $\nabla^{2}\phi(x)v_{x}=\nabla\phi(x)$,
so we can define $\nabla^{2}\phi(x)^{-1}\nabla\phi(x):=v_{x}$. In
the analysis to come, this will be the semantics of the vector $\nabla^{2}\phi(x)^{-1}\nabla\phi(x)$,
and the norm induced by the inverse Hessian should be computed with
this in mind, i.e. $\nabla^{2}f(x)^{-1}$ should be viewed as $\nabla^{2}f(x)^{+}$.
\end{rem*}
Not all convex functions are self-concordant. However, for our purpose,
it suffices to show that we can construct a self-concordant barrier
for any convex set.
\begin{thm}
Any convex set has an $n$-self-concordant barrier.
\end{thm}

Unfortunately, this is an existence result and the barrier function
is expensive to compute. In practice, we construct self-concordant
barriers out of simpler ones:
\begin{lem}
The following functions are self-concordant barriers. We use $\nu$-sc
as a short form for $\nu$-self-concordant barrier.
\begin{itemize}
\item $-\ln x$ is $1$-sc for $\{x\geq0\}$.
\item $-\ln\cos(x)$ is $1$-sc for $\{\left|x\right|\leq\frac{\pi}{2}\}$.
\item $-\ln(t^{2}-\norm x^{2}_{2})$ is $2$-sc for $\{t\geq\norm x_{2}\}$.
\item $-\ln\det X$ is $n$-sc for $\{X\in\R^{n\times n},X\succeq0\}$.
\item $-\ln x-\ln(\ln x+t)$ is $2$-sc for $\{x\geq0,t\geq-\ln x\}$.
\item $-\ln t-\ln(\ln t-x)$ is $2$-sc for $\{t\geq e^{x}\}$.
\item $-\ln x-\ln(t-x\ln x)$ is $2$-sc for $\{x\geq0,t\geq x\ln x\}$.
\item $-2\ln t-\ln(t^{2/p}-x^{2})$ is $4$-sc for $\{t\geq|x|^{p}\}$ for
$p\geq1$.
\item $-\ln x-\ln(t^{p}-x)$ is $2$-sc for $\{t^{p}\geq x\geq0\}$ for
$0<p\leq1$.
\item $-\ln t-\ln(x-t^{-1/p})$ is $2$-sc for $\{x>0,t\geq x^{-p}\}$ for
$p\geq1$.
\item $-\ln x-\ln(t-x^{-p})$ is $2$-sc for $\{x>0,t\geq x^{-p}\}$ for
$p\geq1$.
\end{itemize}
\end{lem}

The following lemma shows how we can combine barriers.
\begin{lem}
If $\phi_{1}$ and $\phi_{2}$ are $\nu_{1}$ and $\nu_{2}$-self-concordant
barriers for $K_{1}$ and $K_{2}$ respectively, then $\phi_{1}+\phi_{2}$
is a $\nu_{1}+\nu_{2}$ self-concordant barrier for $K_{1}\cap K_{2}$.
\end{lem}

\begin{lem}
If $\phi$ is a $\nu$-self-concordant barrier for $K$, then $\phi(Ax+b)$
is a $\nu$-self-concordant barrier for $\{x:Ax+b\in K\}$.
\end{lem}

\begin{xca}
Using the lemmas above, prove that $-\sum^{m}_{i=1}\ln(a^{\top}_{i}x-b_{i})$
is an $m$-self-concordant barrier for the convex set $\{Ax\geq b\}$. 
\end{xca}


\subsection{Convergence of Newton method for self-concordant functions}

Now we are ready to study the convergence of Newton method.
\begin{lem}
\label{lem:Newton_improvement}Given a self-concordant convex function
$f$, consider the iteration 
\[
x'=x-(\nabla^{2}f(x))^{-1}\nabla f(x).
\]
Suppose that $r=\norm{\nabla f(x)}_{\nabla^{2}f(x)^{-1}}<1$, then
we have
\[
\norm{\nabla f(x')}_{\nabla^{2}f(x')^{-1}}\leq\frac{r^{2}}{(1-r)^{2}}.
\]
\end{lem}

\begin{rem}
Note that $\norm{\nabla f(x)}_{\nabla^{2}f(x)^{-1}}=\norm{\nabla^{2}f(x)^{-1}\nabla f(x)}_{x}$
is the step size of the Newton method. This is a measurement of the
error, since the goal is to find $x$ with $\nabla f(x)=0$.
\end{rem}

\begin{proof}
Lemma \ref{lem:self_con_global} shows that
\[
\nabla^{2}f(x')\succeq(1-r)^{2}\nabla^{2}f(x).
\]
and hence
\[
\norm{\nabla f(x')}_{\nabla^{2}f(x')^{-1}}\leq\frac{\norm{\nabla f(x')}_{\nabla^{2}f(x)^{-1}}}{(1-r)}.
\]
To bound $\nabla f(x')$, we calculate that
\begin{align}
\nabla f(x') & =\nabla f(x)+\int^{1}_{0}\nabla^{2}f(x+t(x'-x))(x'-x)dt\nonumber \\
 & =\nabla f(x)-\int^{1}_{0}\nabla^{2}f(x+t(x'-x))(\nabla^{2}f(x))^{-1}\nabla f(x)dt\nonumber \\
 & =\left(\nabla^{2}f(x)-\int^{1}_{0}\nabla^{2}f(x+t(x'-x))dt\right)(\nabla^{2}f(x))^{-1}\nabla f(x).\label{eq:grad_f_sc}
\end{align}
For the term in the bracket, we use Lemma \ref{lem:self_con_global}
to get (note that $\int^{1}_{0}(1-tr)^{2}dt=1-r+\frac{1}{3}r^{2}$):
\[
(1-r+\frac{1}{3}r^{2})\nabla^{2}f(x)\preceq\int^{1}_{0}\nabla^{2}f(x+t(x'-x))dt\preceq\frac{1}{1-r}\nabla^{2}f(x).
\]
Therefore, we have
\[
\norm{(\nabla^{2}f(x))^{-\frac{1}{2}}\left(\nabla^{2}f(x)-\int^{1}_{0}\nabla^{2}f(x+t(x'-x))dt\right)(\nabla^{2}f(x))^{-\frac{1}{2}}}_{\text{op}}\leq\max\left(\frac{r}{1-r},r-\frac{1}{3}r^{2}\right)=\frac{r}{1-r}.
\]
Putting it into (\ref{eq:grad_f_sc}) gives
\begin{align*}
\norm{\nabla f(x')}_{\nabla^{2}f(x)^{-1}} & =\norm{\nabla^{2}f(x)^{-\frac{1}{2}}\nabla f(x')}_{2}\\
 & \le\norm{(\nabla^{2}f(x))^{-\frac{1}{2}}\left(\nabla^{2}f(x)-\int^{1}_{0}\nabla^{2}f(x+t(x'-x))dt\right)(\nabla^{2}f(x))^{-\frac{1}{2}}}_{\text{op}}\norm{(\nabla^{2}f(x))^{-\frac{1}{2}}\nabla f(x)}_{2}\\
 & \leq\frac{r}{1-r}\cdot r\\
 & =\frac{r^{2}}{1-r}.
\end{align*}
\end{proof}
Finally, we bound the error of the current iterate in terms of $\|\nabla f(x)\|_{\nabla^{2}f(x)^{-1}}$.
\begin{lem}
\label{lem:self_concordant_err}Given $x$ such that $\|\nabla f(x)\|_{\nabla^{2}f(x)^{-1}}\leq\frac{1}{6}$,
we have that
\begin{itemize}
\item $\|x-x^{*}\|_{x^{*}}\leq2\|\nabla f(x)\|_{\nabla^{2}f(x)^{-1}}$,
\item $\|x-x^{*}\|_{x}\leq\frac{4}{3}\|\nabla f(x)\|_{\nabla^{2}f(x)^{-1}}$,
\item $f(x)\leq f(x^{*})+2\|\nabla f(x)\|^{2}_{\nabla^{2}f(x)^{-1}}.$
\end{itemize}
\end{lem}

\begin{proof}
Let $r=\|x-x^{*}\|_{x}$. Suppose that $r\leq\frac{1}{4}$. Note that
\[
\nabla f(x)=\nabla f(x)-\nabla f(x^{*})=\int^{1}_{0}\nabla^{2}f(x^{*}+t(x-x^{*}))(x-x^{*})dt.
\]
Using that $\nabla^{2}f(x^{*}+t(x-x^{*}))\succeq(1-(1-t)r)^{2}\nabla^{2}f(x)$
(Lemma \ref{lem:self_con_global}), we have
\begin{align*}
\|\nabla f(x)\|_{\nabla^{2}f(x)^{-1}} & =\norm{\int^{1}_{0}\nabla^{2}f(x^{*}+t(x-x^{*}))(x-x^{*})\,dt}_{\nabla^{2}f(x)^{-1}}\\
 & \geq\int^{1}_{0}(1-(1-t)r)^{2}\|x-x^{*}\|_{x}dt\\
 & =\left(1-r+\frac{r^{2}}{3}\right)r\geq\frac{3r}{4}.
\end{align*}

Using $\|\nabla f(x)\|_{\nabla^{2}f(x)^{-1}}\leq\frac{1}{6}$, we
have indeed $r\leq\frac{1}{4}$ (our lower bound above is a non-decreasing
function). Using Lemma \ref{lem:self_con_global} again, we have
\[
\|x-x^{*}\|_{x^{*}}\leq\frac{\|x-x^{*}\|_{x}}{1-r}\leq\frac{4}{3}\frac{1}{1-\frac{1}{4}}\|\nabla f(x)\|_{\nabla^{2}f(x)^{-1}}\leq2\|\nabla f(x)\|_{\nabla^{2}f(x)^{-1}}.
\]

For the bound for $f(x)$, we have that
\[
f(x)=f(x^{*})+\left\langle \nabla f(x^{*}),x-x^{*}\right\rangle +\int^{1}_{0}(1-t)(x-x^{*})^{\top}\nabla^{2}f(x^{*}+t(x-x^{*}))(x-x^{*})dt.
\]
Using that $\nabla^{2}f(x^{*}+t(x-x^{*}))\preceq\frac{1}{(1-(1-t)r)^{2}}\nabla^{2}f(x)\preceq\frac{16}{9}\nabla^{2}f(x)$,
we have
\begin{align*}
f(x) & \leq f(x^{*})+\frac{16}{9}\int^{1}_{0}(1-t)dt\cdot\|x-x^{*}\|^{2}_{x}\\
 & \leq f(x^{*})+\|x-x^{*}\|^{2}_{x}\\
 & \le f(x^{*})+2\norm{\nabla f(x)}^{2}_{\nabla^{2}f(x)^{-1}}.
\end{align*}
\end{proof}

\subsection{Main Algorithm and Analysis}

\begin{algorithm2e}[H]

\caption{$\mathtt{InteriorPointMethod}$}

\SetAlgoLined

\textbf{Input:} A $\nu$-self-concordant barrier $\phi$ for $K$,
the minimizer $x$ of $\phi$. 

Define $f_{t}(x)=tc^{\top}x+\phi(x)$. $t=\frac{1}{6}\|c\|^{-1}_{\nabla^{2}\phi(x)^{-1}}$.

\While{$t\leq\frac{\nu+\sqrt{\nu}}{\epsilon}$}{

$x\leftarrow x-\nabla^{2}f_{t}(x)^{+}\nabla f_{t}(x)$.

$t\leftarrow(1+h)t$ with $h=\frac{1}{9\sqrt{\nu}}$.

}

\Return $x$.

\end{algorithm2e}

We first explain the initialization and termination conditions. We
start with the point $x^{(0)}=\arg\min\phi(x)$ (or a point very close
to it). The choice of initial $t$ is to ensure that the norm of the
initial gradient in the inverse Hessian norm is bounded by a convenient
constant less than $1$ (we use $1/6$). 

For the termination, intuitively, we can see that $\min f_{t}(x_{t})$
tends to optimality as $t\rightarrow\infty$. We first need a lemma
showing that the gradient of $\phi$ is small.
\begin{lem}[Duality Gap]
Suppose that $\phi$ is a $\nu$-self-concordant barrier. For any
$x,y\in K$, we have that
\[
\left\langle \nabla\phi(x),y-x\right\rangle \leq\nu.
\]
\end{lem}

\begin{proof}
Let $\alpha(t)=\left\langle \nabla\phi(z_{t}),y-x\right\rangle $
where $z_{t}=x+t(y-x)$. Then, we have
\[
\alpha'(t)=\left\langle \nabla^{2}\phi(z_{t})(y-x),y-x\right\rangle .
\]
Note that
\[
\alpha(t)\leq\norm{\nabla\phi(z_{t})}_{\nabla^{2}\phi(z_{t})^{-1}}\norm{y-x}_{\nabla^{2}\phi(z_{t})}\leq\sqrt{\nu}\norm{y-x}_{\nabla^{2}\phi(z_{t})}.
\]
Hence, we have $\alpha'(t)\geq\frac{1}{\nu}\alpha(t)^{2}.$ If $\alpha(0)\leq0$,
then we are done. Otherwise, $\alpha$ is increasing and hence, $\alpha(1)>0$.
Therefore, 
\[
\frac{1}{\alpha(1)}\leq\frac{1}{\alpha(0)}-\frac{1}{\nu}
\]
and we have $\alpha(0)\leq\nu$.
\end{proof}
\begin{lem}[Duality Gap]
 Suppose that $\phi$ is a $\nu$-self-concordant barrier, we have
that
\[
\left\langle c,x_{t}\right\rangle \leq\left\langle c,x^{*}\right\rangle +\frac{\nu}{t},
\]
where $x_{t}=\argmin f_{t}(x)$. More generally, for any $x$ such
that $\|tc+\nabla\phi(x)\|_{(\nabla^{2}\phi(x))^{-1}}\leq\frac{1}{6}$,
we have that
\[
\left\langle c,x\right\rangle \leq\left\langle c,x^{*}\right\rangle +\frac{\nu+\sqrt{\nu}}{t}.
\]
\end{lem}

\begin{proof}
Let $x^{*}$ be a minimizer of $c^{\top}x$ on $K$. By optimality
of $x_{t}$ for $f_{t}$, we have $tc+\nabla\phi(x_{t})=0$. Therefore,
we have 
\[
\left\langle c,x_{t}\right\rangle -\left\langle c,x^{*}\right\rangle =\frac{1}{t}\left\langle \nabla\phi(x_{t}),x^{*}-x_{t}\right\rangle \leq\frac{\nu}{t}.
\]

For the second result, let $f(x)=tc^{\top}x+\phi(x)$. Lemma \ref{lem:self_concordant_err}
shows that
\[
\|x-x_{t}\|_{x}\leq2\|\nabla f(x)\|_{\nabla^{2}f(x)^{-1}}\leq\frac{1}{3}.
\]
Hence,
\begin{align*}
\left\langle c,x-x_{t}\right\rangle  & \leq\|c\|_{\nabla^{2}\phi(x)^{-1}}\|x-x_{t}\|_{x}\leq\frac{1}{3}\|c\|_{\nabla^{2}\phi(x)^{-1}}
\end{align*}
Using $c=\frac{tc+\nabla\phi(x)}{t}-\frac{\nabla\phi(x)}{t}$, we
have
\begin{align*}
\left\langle c,x-x_{t}\right\rangle  & \leq\frac{1}{3t}\left(\|tc+\nabla\phi(x)\|_{\nabla^{2}\phi(x)^{-1}}+\|\nabla\phi(x)\|_{\nabla^{2}\phi(x)^{-1}}\right)\\
 & \leq\frac{1}{3t}\left(\frac{1}{6}+\sqrt{\nu}\right)\leq\frac{\sqrt{\nu}}{t}.
\end{align*}
This gives the result.
\end{proof}
Hence, it suffices to end with $t=(\nu+\sqrt{\nu})/\varepsilon$,
which is exactly the same as in the linear-programming case above.
\begin{thm}
Given a $\nu$-self-concordant barrier $\phi$ and its minimizer,
we can find $x\in K$ such that $c^{\top}x\leq c^{\top}x^{*}+\epsilon$
in 
\[
O(\sqrt{\nu}\log(\frac{\nu}{\epsilon}\|c\|_{\nabla^{2}\phi(x^{(0)})^{-1}}))
\]
iterations.
\end{thm}

\begin{proof}
We prove by induction that $\|\nabla f_{t}(x)\|_{\nabla^{2}f_{t}(x)^{-1}}\leq\frac{1}{6}$
at the beginning of each iteration. This is true at the beginning
by the definition of initial $t$. By Lemma \ref{lem:Newton_improvement},
after the Newton step, we have
\[
\|\nabla f_{t}(x)\|_{\nabla^{2}f_{t}(x)^{-1}}\leq\left(\frac{1/6}{1-1/6}\right){}^{2}=\frac{1}{25}.
\]
Let $t'=(1+h)t$ with $h=\frac{1}{9\sqrt{\nu}}$. Note that $\nabla f_{t'}(x)=(1+h)tc+\nabla\phi(x)$
and hence $\nabla f_{t'}(x)=(1+h)\nabla f_{t}(x)-h\nabla\phi(x)$
and $\nabla^{2}f_{t'}(x)=\nabla^{2}f_{t}(x)=\nabla^{2}\phi(x).$ Therefore,
we have that
\begin{align*}
\|\nabla f_{t'}(x)\|_{\nabla^{2}f_{t'}(x)^{-1}} & =\|(1+h)\nabla f_{t}(x)-h\nabla\phi(x)\|_{\nabla^{2}f_{t}(x)^{-1}}\\
 & \leq(1+h)\|\nabla f_{t}(x)\|_{\nabla^{2}f_{t}(x)^{-1}}+h\|\nabla\phi(x)\|_{\nabla^{2}\phi(x)^{-1}}\\
 & \leq\frac{1+h}{25}+h\sqrt{\nu}\leq\frac{1}{6}.
\end{align*}

This completes the induction.
\end{proof}

